%% kapa: exam=1 round=32 year=2021 period=1 form=A subject=civil_law subject_ko=민법 seq=8 number=8
%% source: Q-Net 감정평가사 기출문제 (cst00309) artlSeq=5208397 file=extracted/1008/5208397/A형/2021년도 제32회 감정평가사 1차 1교시 A형.pdf
%% conversion: kordoc 4.24.51 -> tools/gen1.py (auto); review=proofread-sonnet; flags=none
\begin{question}{8}
물건에 관한 설명으로 옳지 않은 것은? (다툼이 있으면 판례에 따름)
\begin{choices}
  \item 주물 소유자의 사용에 공여되는 물건이라도 주물 자체의 효용과 직접 관계가 없으면 종물이 아니다.
  \item 「입목에 관한 법률」에 의하여 소유권보존등기를 한 수목의 집단은 저당권의 객체가 된다.
  \item 종물과 주물의 관계에 관한 법리는 권리 상호간에도 적용될 수 있다.
  \item 분필절차를 거치지 않은 1필의 토지의 일부에 대해서도 저당권을 설정할 수 있다.
  \item 저당권의 효력은 저당부동산의 종물에 미치므로 경매를 통하여 저당부동산의 소유권을 취득한 자는 특별한 사정이 없는 한 종물의 소유권을 취득한다.
\end{choices}
\end{question}
\answer{4}
